\chapter{Sharp Runtime: Strings, Streams, and a Real Audit History}
\label{ch:sharp-runtime-namespaces}

This chapter covers two smaller, practical corners of \texttt{sharp-runtime}'s API --- its
text and stream types --- and then the part of its history most worth learning from: what
happened when the project's own completed task tracker was deliberately \emph{not} trusted as
proof of correctness, and independent audits went looking for what it might have missed.

\section{String and StringBuilder: a static-only utility versus a real mutable buffer}

\cnaclass{System::String} is deliberately \emph{not} a wrapper class at all --- its default
constructor and destructor are both explicitly deleted, making it a pure static-method utility
operating on ordinary \texttt{std::string} values (\texttt{IsNullOrEmpty}, \texttt{Compare},
the \texttt{IndexOf}/\texttt{LastIndexOf} family, \texttt{Split}, \texttt{Substring},
\texttt{Trim*}, \texttt{PadLeft}/\texttt{PadRight}). \cnaclass{Text::StringBuilder}, by
contrast, is a real, stateful, mutable buffer type --- the many-overloaded \texttt{Append}
(string, char, int, double, float, bool, and the \texttt{longcs} alias from
Chapter~\ref{ch:design-philosophy} in Volume~I), \texttt{AppendLine}, \texttt{Insert},
\texttt{Remove}, \texttt{Replace}, and a nested \texttt{ChunkEnumerator} class mirroring real
.NET's own chunked-buffer enumeration shape.

\section{Stream: a lightweight subset, honest about it}

\cnaclass{IO::Stream}'s own class comment states its scope plainly: ``lightweight subset of
the .NET Stream API.'' \texttt{Read(buffer, offset, count)}, \texttt{Close()}, and
\texttt{getLengthProperty()} are pure virtual; \texttt{Write}/\texttt{WriteByte} default to
throwing \texttt{NotSupportedException} unless a writable stream subclass overrides them ---
the same shape real .NET's own read-only stream implementations follow, rather than forcing
every stream subclass to implement a full read/write/seek contract regardless of whether it
actually supports all three.

\section{The post-stabilization audit: distrusting a complete-looking task list on purpose}

\texttt{sharp-runtime}'s own ticket backlog --- 1{,}709 tickets across sixteen categories ---
reached 100\% \texttt{done} at one point in the project's history. What happened next is the
single most instructive episode in this project's own story, and worth recounting in full,
because it is the same discipline Chapter~\ref{ch:project-practice} named as this whole
ecosystem's most distinctive trait, applied at real scale: rather than treating a fully
checked-off ticket table as proof the runtime was correct, the project commissioned a fresh,
independent audit specifically built to distrust that conclusion. Seven parallel, read-only
audit agents were dispatched across eight risk categories --- API inconsistencies,
undocumented deviations from real .NET, remaining \texttt{NotImplementedException}s, partial
or stub classes, platform-specific issues, exception-type mismatches, silently-wrong
behavior, and missing high-risk test coverage --- each one required to verify every finding
directly against real .NET source and against the project's own documented, accepted
deviations list before it could be reported at all, specifically to filter out
plausible-sounding but wrong findings before they ever reached a human.

The result was twenty concrete findings --- nine high-severity, eight medium, three low ---
filed as new tickets and subsequently fixed. One is worth naming in full, because of how
quietly it could have shipped: \texttt{Convert::ToXxx(double/float)} was implemented as a
bare \texttt{static\_cast<T>()}, meaning simple truncation rather than round-half-to-even
rounding. A direct repro confirmed \texttt{Convert::ToInt32(2.9)} returned \texttt{2} where
real .NET returns \texttt{3}, and \texttt{Convert::ToInt32(3.5)} returned \texttt{3} where real
.NET returns \texttt{4} --- and, more strikingly, \texttt{ToUInt32(double)}'s own doc comment
already said \texttt{"(truncates)"}, as though the wrong behavior had been documented as
intentional rather than caught as a bug.

\section{A second, independent audit, and what sanitizers found that reasoning alone did not}

A further, separate audit produced seven more findings, six fixed in the same session ---
among them a stream-backed ZIP archive implementation that never actually wrote its output
back to the underlying stream in create or update mode, and five of this project's own
destructors (spanning compression streams and XML writing) that could call
\texttt{std::terminate()} if an I/O failure occurred during RAII cleanup, since a
C\texttt{++} destructor that lets an exception escape terminates the program outright rather
than propagating the failure the way a .NET \texttt{Dispose()} call safely could.

A first-ever full run of the full sanitizer trio --- ThreadSanitizer, AddressSanitizer, and
UndefinedBehaviorSanitizer --- found real bugs no amount of code review alone had caught:
ThreadSanitizer found one genuine production deadlock (a \texttt{Channel} implementation
missing a \texttt{notify\_all} call, a lost-wakeup bug affecting multiple blocked waiters)
alongside three test-only races; AddressSanitizer found one confirmed heap-buffer-overflow (an
aligned-reallocation helper copying the \emph{old}, smaller allocation's size into a newly
enlarged buffer) plus five real memory leaks in XML document node creation; and
UndefinedBehaviorSanitizer found one genuine undefined-behavior call reaching into vendored
\texttt{miniz} code, from an empty vector's \texttt{.data()} pointer reaching a
non-null-parameter-annotated function.

\section{Why this history belongs in a book, not just a changelog}

None of these bugs are individually remarkable --- a truncation-instead-of-rounding bug, a
missed \texttt{notify\_all}, a leaking XML method. What is worth taking from this chapter is
the project's own response to having a fully completed task tracker: rather than reading
``100\% done'' as ``verified correct,'' it treated that state as precisely the moment a fresh,
independently-skeptical audit was most valuable, and built that audit specifically to check
claims against real .NET source rather than against the project's own prior conclusions. This
is the same discipline \texttt{xna4-spec} auditing (Chapter~\ref{ch:xna4-spec-auditing})
brings to CNA's own XNA surface, applied here one layer further down, to the foundation layer
everything else in this book's Volume~I ultimately stands on.
